%----------
%	GLOSARIO DE TÉRMINOS
%----------

Esta memoria está redactada en español, pero emplea un número apreciable de
términos en inglés.

Se traducen al español los términos que cuentan con un equivalente asentado en la
disciplina correspondiente. Es el caso de \textit{codebook}, que aquí aparece
siempre como libro de códigos por ser la forma habitual en análisis de contenido.

Se mantienen en inglés tres tipos de término. El primero son los que carecen de un
equivalente consolidado en español, de modo que traducirlos introduciría ambigüedad
donde el original es unívoco, como ocurre con \textit{prompt} o \textit{token}. El
segundo son los nombres propios de una tecnología, una arquitectura o una
especificación concreta, como \textit{Transformer} o \textit{Agent Skills}, que
funcionan como denominación oficial y cuya traducción dificultaría reconocerlos. El
tercero son las expresiones que identifican un régimen experimental reconocible en
la literatura, como \textit{zero-shot}, cuya versión española obligaría a una
perífrasis en cada aparición.

En todos los casos, la primera vez que un término aparece en el texto se ofrece su
equivalente en español junto a la forma inglesa. Las tablas siguientes
recogen los términos y acrónimos empleados, con su equivalencia y una breve nota
sobre el criterio aplicado.

Conviene justificar una segunda decisión lingüística, esta vez de género
gramatical. En inglés, \textit{agent} carece de género. En español, «agente» es un
sustantivo común en cuanto al género, de modo que admite tanto el masculino como el
femenino sin alterar su significado~\cite{rae_nueva_gramatica}. Esta memoria emplea
deliberadamente el femenino, «la agente», por coherencia con su propio objeto de
estudio. Un trabajo que analiza el masculino genérico como una de sus variables de
análisis (V26) y que deriva sus criterios de guías de lenguaje no sexista
difícilmente puede recurrir al masculino por defecto en su redacción cuando la lengua
ofrece una alternativa igualmente correcta. La decisión sigue el criterio que esas
mismas guías proponen, el de favorecer «formulaciones que den lugar a
representaciones menos masculinizadas de las realidades a las que se
refieren»~\cite{yufera2023}.

\begin{table}[H]
    \centering
    \begingroup
    \scriptsize
    \setlength{\tabcolsep}{4pt}
    \ttabbox[\FBwidth]{
        \caption{Acrónimos en inglés empleados en la memoria}
        \label{tab:glosario-acronimos}
    }{
        \begin{tabularx}{\textwidth}{
            >{\hsize=0.45\hsize\raggedright\arraybackslash}X
            >{\hsize=1.25\hsize\raggedright\arraybackslash}X
            >{\hsize=1.30\hsize\raggedright\arraybackslash}X
        }
            \toprule
            \textbf{Sigla} & \textbf{Forma desarrollada} & \textbf{Equivalente en español} \\
            \midrule
            API & \textit{Application Programming Interface} & Interfaz de programación de aplicaciones \\
            HCI & \textit{Human-Computer Interaction} & Interacción persona-computador \\
            LLM & \textit{Large Language Model} & Gran modelo de lenguaje \\
            PLE & \textit{Per-Layer Embeddings} & Representaciones vectoriales por capa \\
            RAG & \textit{Retrieval-Augmented Generation} & Generación aumentada por recuperación \\
            TDM & \textit{Text and Data Mining} & Minería de textos y datos \\
            TF-IDF & \textit{Term Frequency-Inverse Document Frequency} & Frecuencia de término y frecuencia inversa de documento \\
            \bottomrule
        \end{tabularx}
    }
    \endgroup
\end{table}



\begin{table}[H]
    \centering
    \begingroup
    \scriptsize
    \setlength{\tabcolsep}{4pt}
    \ttabbox[\FBwidth]{
        \caption{Glosario de términos en inglés empleados en la memoria}
        \label{tab:glosario-terminos}
    }{
        \begin{tabularx}{\textwidth}{
            >{\hsize=0.85\hsize\raggedright\arraybackslash}X
            >{\hsize=0.95\hsize\raggedright\arraybackslash}X
            >{\hsize=1.20\hsize\raggedright\arraybackslash}X
        }
            \toprule
            \textbf{Término} & \textbf{Equivalente en español} & \textbf{Nota} \\
            \midrule
            \textit{Agent Skills} & Habilidades agénticas & Denominación de la especificación que empaqueta conocimiento experto en ficheros cargables bajo demanda. Se conserva el original \\
            \textit{backend} & Parte servidora & Componente que atiende las peticiones del sistema \\
            \textit{baseline} & Condición de control & Configuración de referencia contra la que se compara la propuesta \\
            \textit{benchmark} & Banco de pruebas & Conjunto estandarizado de tareas para comparar sistemas \\
            \textit{circuit breaker} & Mecanismo de corte & Salvaguarda que aborta una ejecución cuando se acumulan fallos consecutivos \\
            \textit{codebook} & Libro de códigos & \textbf{Se traduce} en toda la memoria, por ser la forma asentada en análisis de contenido \\
            \textit{embedding} & Representación vectorial & Vector que codifica un texto. Se conserva el original por su uso mayoritario \\
            \textit{few-shot} & Con pocos ejemplos & Régimen en que la instrucción incluye algunos casos ya resueltos \\
            \textit{fine-tuning} & Ajuste fino & Reentrenamiento de un modelo sobre datos específicos de la tarea \\
            \textit{ground truth} & Anotación de referencia & Etiquetas expertas que sirven de patrón de comparación \\
            \textit{hardware} & --- & Voz admitida por la Real Academia Española, se emplea sin traducir \\
            \textit{human-in-the-loop} & Persona en el bucle & Paradigma en que el sistema asiste y la persona conserva la decisión \\
            \textit{mainstream} & Corriente principal & Referido a los medios de comunicación de mayor difusión \\
            \textit{progressive disclosure} & Revelación progresiva & Cargar en cada momento solo la información necesaria para la tarea \\
            \textit{prompt} & Instrucción & Texto que recibe el modelo para resolver la tarea \\
            \textit{prompting} & Ingeniería de instrucciones & Diseño y formulación de las instrucciones que recibe el modelo \\
            \textit{shard} & Fragmento & Porción del corpus que se procesa en paralelo \\
            \textit{skill} & Habilidad & Unidad de conocimiento cargable bajo demanda por la agente \\
            \textit{software} & --- & Voz admitida por la Real Academia Española, se emplea sin traducir \\
            \textit{token} & Unidad léxica & Unidad mínima de texto que procesa el modelo. Se conserva por precisión técnica \\
            \textit{tool-calling} & Llamada a herramientas & Mecanismo nativo de algunos proveedores para invocar funciones \\
            \textit{toolkit} & Conjunto de herramientas & Paquete de utilidades orientado a una tarea concreta \\
            \textit{Transformer} & --- & Arquitectura de red neuronal. Nombre propio, no se traduce \\
            \textit{zero-shot} & Sin ejemplos previos & Régimen en que el modelo resuelve la tarea sin casos resueltos en la instrucción \\
            \bottomrule
        \end{tabularx}
    }
    \endgroup
\end{table}

